The enumeration of fixed-margin arrays has long been studied in
statistics and combinatorics; Diaconis and Gangolli~\cite{DG1995}
survey its applications and exact and approximate methods.
Ordinary contingency tables are a special case of
\eqref{eq:feasible-set}, since the margins imply
$X_{ij}\le\min\{r_i,c_j\}$. Even their exact enumeration is difficult:
Dyer, Kannan, and Mount~\cite[Theorem~1]{DKM1997} proved
$\#\mathrm P$-completeness for tables with only two rows.
Approximate counting can nevertheless be efficient under much broader
conditions than exact counting.

For ordinary tables, Dyer, Kannan, and Mount~\cite{DKM1997}
developed geometric sampling and counting methods for sufficiently
large margins. Morris~\cite{Morris2002} weakened the required margin
bounds by refining the expanded-polytope and rounding argument.
Dyer and Greenhill~\cite{DG2000} gave polynomial-time approximate
counting and sampling for arbitrary two-row margins. Cryan and
Dyer~\cite{CD2003} obtained an FPRAS for any fixed number of rows by
combining dynamic programming and volume estimation; Dyer~\cite{Dyer2003}
subsequently improved the fixed-row algorithm using dynamic programming.
Gopalan, Klivans, Meka, \v{S}tefankovi\v{c}, Vempala, and
Vigoda~\cite[Theorem~1.3]{GKM2011} later gave a deterministic FPTAS
for a fixed number of rows.
These results already permit binary-encoded margins. Their restrictions
concern the number of rows or the sizes of the margins, rather than a
unary encoding convention.

Individual caps connect the problem to several other counting models.
When all caps are zero or one, a table is a subgraph of the allowed
bipartite graph with prescribed vertex degrees. Jerrum, Sinclair, and
Vigoda~\cite[Corollary~8.1]{JSV2004} give an FPRAS for this entire
special case as a consequence of their permanent algorithm; all-one
margins give perfect matchings. Bez\'akov\'a, Bhatnagar, and
Vigoda~\cite{BBV2007} developed a direct annealing algorithm for
binary tables, also allowing arbitrary forbidden cells.
More generally, capped tables are
integer flows in a bipartite network. Cryan, Dyer, and
Randall~\cite{CDR2010} gave an FPRAS for integral flows with sufficiently
large tight capacities, and for cell-bounded tables with a fixed number
of rows and arbitrary positive margins and caps. They identified
unrestricted approximation as an open problem in their 2010 paper.
Their tight-capacity reduction and maximum-capacity spanning-tree
representation are direct antecedents of our large-cell construction.
Here a forest and soft constraints accommodate arbitrary mixtures of
small and large bounds, with both dimensions variable and zero bounds
permitted.

Convex optimization gives a complementary family of estimates.
Barvinok~\cite[Theorem~1.3]{Barvinok2009} bounds ordinary weighted
table counts, including forbidden cells, within a factor
$N^{O(m+n)}$, where $N$ is the total margin. Barvinok, Luria,
Samorodnitsky, and Yong~\cite{BLSY2010} obtain randomized relative
approximations with complexity $N^{O(\log N)}$ for specified
smooth-margin classes. This scale is quasipolynomial in the numerical
total, rather than in its binary length. Br\"and\'en, Leake, and
Pak~\cite[Theorem~2.1]{BLP2021} give capacity bounds for arbitrary
individual caps, with explicit multiplicative factors depending on
the margins and dimensions. Their bounds extend earlier estimates
using Lorentzian polynomials. A subsequent preprint of September 30,
2026, by Leake and Mohammadi Yekta~\cite[Theorems~1.1 and~1.3]{LMY2026}
extends capacity-based estimates to finite nonnegative integer fibers
of totally unimodular systems, retaining an explicit multiplicative
factor depending on the system's dimensions and right-hand side.
Theorem~\ref{thm:main} supplies
any requested relative tolerance in polynomial binary time, without
a smoothness assumption.

The geometric problem here is to count all integer points of a
transportation polytope. These need not all be vertices: the
$2\times2$ tables with every margin equal to two form a segment
with three integer points and two vertices. Guo and
Jerrum~\cite[Proposition~4.2]{GJ2023} give an FPRAS for the vertices
of promised $0/1$ polytopes $\{x:0\le Ax\le1\}$ defined by network
matrices; this unit-bound vertex result has a different scope.

Counting also provides approximate sampling through the
self-reduction principle of Jerrum, Valiant, and
Vazirani~\cite[Theorem~6.3]{JVV1986}. Section~\ref{sec:flows}
implements this principle by bisecting cell intervals, with exact
feasibility tests at every branch. For an explicitly given directed
multigraph with finite integer arc bounds and prescribed vertex
balances, it yields both an FPRAS for integer arc-value vectors and
a sampler that always returns a feasible vector when one exists.
The latter has total-variation error at most $\tau$ and an
all-execution cost polynomial in the input length and $\tau^{-1}$.
A separate companion~\cite[Theorem~1.1]{OpenAIExactTables2026}
proves exact uniform sampling in polynomial expected bit time for
ordinary tables with arbitrary margins and no input cell bounds or
forbidden positions. Its bounded approximate sampler has cost
polynomial in the input length and the number of accuracy bits. Those stronger sampling
guarantees concern the uncapped setting; the flow consequence here
has the total-variation and cost bounds just stated.
